\originalchapter{闭性、分离与共轭}\label{chap:3}
本章对应原第6--8讲的闭集交、分离与共轭部分, 第64--103页. 三个看似不同的问题, 即最优解是否存在、线性像是否闭、消去变量后的值函数是否闭, 都可转化为嵌套闭集的交是否非空.

\originalsection{嵌套闭集与回缩性}
若 $C_{k+1}\subseteq C_k$ 都非空闭, 且某个 $C_k$ 紧, 则交集非空: 在每个后续集合选点, 紧性给收敛子序列, 嵌套性使极限属于每个固定集合. 不紧时这一结论失败, 如 $C_k=[k,\infty)$.

\begin{definition}
对嵌套非空闭凸集 $C_k$, 若 $x_k\in C_k$, $\norm{x_k}\to\infty$ 且 $x_k/\norm{x_k}\to d/\norm d$, $d\ne0$, 称其为渐近序列. 若最终 $x_k-d\in C_k$, 称这一渐近序列可回缩. 若每个渐近序列都可回缩, 称集合序列可回缩. 单个闭凸集的可回缩性由常值集合序列定义.
\end{definition}
任何无界的选择序列都存在渐近子序列. 其方向属于每个 $\rec C_j$: 固定 $j$, 后续点都在 $C_j$, 对规范化线段取极限即可使用衰退锥的证明. 回缩性进一步要求在已很远的可行点上向反方向移动固定距离仍可行, 它比“正方向可以无限移动”更强.

\begin{lemma}
闭半空间、多面体是可回缩集. 有限个可回缩集合序列的逐项交及笛卡尔积仍可回缩. 凸紧集与可回缩凸集的向量和可回缩.
\end{lemma}
\begin{proof}
半空间记为 $a\T x\le b$. 对渐近方向 $d$, 必有 $a\T d\le0$. 若等于0, 平移 $-d$ 不改变不等式; 若小于0, 则 $a\T x_k\to-\infty$, 最终也允许平移. 多面体只有有限个条件, 逐一应用即可. 交集的渐近序列是每个组成序列的渐近序列, 故可以同时回缩.

对积, 各分量中方向非零的项按相应增长比例作渐近序列, 使用该分量的回缩性; 方向为零的项无须移动. 对和, 写 $x_k=u_k+v_k$, 紧集分量 $u_k$ 有界. 则 $v_k$ 具有与 $x_k$ 相同的渐近方向, 最终 $v_k-d$ 可行, 所以 $x_k-d=u_k+(v_k-d)$ 可行.
\end{proof}
\begin{theorem}\label{thm:retractive-intersection}
可回缩的嵌套非空闭凸集合序列具有非空交集. 更一般地, 设 $X$ 为可回缩闭凸集, 每个 $X\cap C_k$ 非空, 并令 $R=\bigcap_k\rec C_k$, $L=\bigcap_kL_{C_k}$. 若 $\rec X\cap R\subseteq L$, 则 $\bigcap_k(X\cap C_k)$ 非空.
\end{theorem}
\begin{proof}
在每个 $C_k$ 中取离原点最近的点 $p_k$. 若此序列有界, 收敛子序列的极限属于交集. 若无界, 取渐近子序列, 方向为 $d$. 回缩性给出最终 $p_k-d\in C_k$, 但
$\norm{p_k-d}^2=\norm{p_k}^2-2p_k\T d+\norm d^2<\norm{p_k}^2$, 因为 $p_k\T d\to+\infty$. 这与最近点定义矛盾.

对一般情况, 渐近方向 $d$ 属于 $\rec X\cap R$, 因而属于每个 $L_{C_k}$. 所以 $x_k-d\in C_k$ 始终成立; 而 $X$ 的回缩性保证它最终属于 $X$. 于是 $X\cap C_k$ 序列可回缩, 用第一部分得结论.
\end{proof}
$X$ 的回缩性不能删掉. 取闭凸集 $X=\{(x,y):x>0,y\ge1/x\}$ 与 $C_k=\{(x,y):y\le1/k\}$, 各个交非空, 全部交却为空. $\rec X\cap\bigcap_k\rec C_k$ 是非负水平射线, 包含于每个 $C_k$ 的线性空间, 但 $X$ 不可回缩: 点 $(k,1/k)$ 向左平移1后不再属于 $X$.

\begin{corollary}\label{cor:qp-attainment}
在非空多面体上, 半正定二次函数只要下确界有限, 就取到最小值. 线性规划是其特例.
\end{corollary}
\begin{proof}
写 $f(x)=x\T Qx+c\T x$, $Q\succeq0$, 取 $\gamma_k\downarrow f^*$, $\gamma_k>f^*$. 下水平集的共同衰退锥为 $R=\{d:Qd=0,c\T d\le0\}$, 共同线性空间为 $L=\{d:Qd=0,c\T d=0\}$. 若 $d\in\rec X\cap R$ 且 $c\T d<0$, 任一可行点沿 $d$ 使目标趋于 $-\infty$, 与有限下确界矛盾. 因而 $\rec X\cap R\subseteq L$. 多面体可回缩, 上一定理使嵌套可行水平集的交非空, 交中点恰为最优点.
\end{proof}

\originalsection{线性像与部分极小化}
\begin{theorem}\label{thm:closed-image}
设 $C$ 非空闭凸, $A$ 线性. 若 $\rec C\cap\ker A\subseteq L_C$, 则 $AC$ 闭. 若 $X$ 是可回缩闭凸集, 则条件 $\rec X\cap\rec C\cap\ker A\subseteq L_C$ 保证 $A(X\cap C)$ 闭. 特别地, 多面体的线性像闭.
\end{theorem}
\begin{proof}
取 $y_k=Ax_k\to y$. 令 $\varepsilon_k\downarrow0$ 且选择子序列使 $\norm{y_k-y}\le\varepsilon_k$, 再令 $D_k=\{x\in C:\norm{Ax-y}\le\varepsilon_k\}$. 各 $D_k$ 非空闭凸, 共同衰退锥为 $\rec C\cap\ker A$, 线性空间包含 $L_C\cap\ker A$. 条件及定理\ref{thm:retractive-intersection}, 取 $X=\R^n$, 保证交非空; 交中点满足 $Ax=y$.

第二种情况在同一 $D_k$ 外再交 $X$, 并用一般回缩交定理. 多面体的线性像可令 $C=\R^n$, 此时 $L_C=\R^n$, 条件自动成立.
\end{proof}
原第1讲第13页的例子也可精确写出: $C_1=\{(u,v):u>0,v>0,uv\ge1\}$ 与 $C_2=\{(0,v):v\in\R\}$ 都闭凸, 但 $C_1+C_2=\{(u,v):u>0,v\in\R\}$ 不闭. 前者闭是因为趋于 $u=0$ 时 $v$ 必趋无穷, 不产生有限漏点; 与竖直线相加后这个阻止有限极限的条件消失. 同样, 将 $C_1$ 投影到第一坐标只得到开半轴.

最简单的充分条件是 $\rec C\cap\ker A=\{0\}$, 这时预像截集本身紧. 向量和是积在线性求和映射下的像; 因而若 $d_i\in\rec C_i$ 且 $\sum_i d_i=0$ 只能在各 $d_i=0$ 时成立, 则 $\sum_iC_i$ 闭. 特别地, 一个闭凸集与一个凸紧集的和闭, 以及 $\rec C_1\cap\rec C_2=\{0\}$ 时 $C_1-C_2$ 闭.

\begin{theorem}\label{thm:partial-min}
设 $F:\R^{n+m}\to(-\infty,+\infty]$ 闭真凸, $p(x)=\inf_zF(x,z)$. 则 $p$ 凸. 若某个纤维水平集 $\{z:F(x_0,z)\le\gamma\}$ 非空紧, 则 $p$ 闭且真, 且对每个 $x\in\dom p$, 纤维极小值都取到, 极小点集合非空紧.
\end{theorem}
\begin{proof}
对两个近似极小点作凸组合, 目标误差可任意小, 得 $p$ 凸. 记 $P(x,z,w)=(x,w)$. 总有 $P(\epi F)\subseteq\epi p\subseteq\cl(P(\epi F))$: 第二个包含由在固定 $x$ 处取任意接近下确界的 $z$ 得到.

某个纤维水平集紧, 说明上图没有非零水平纤维衰退方向 $(0,d,0)$. 否则该方向会使这一个紧集包含非零射线. 因而 $\rec(\epi F)\cap\ker P=\{0\}$, 线性像定理使 $P(\epi F)$ 闭, 两个包含变为等号. 对任一非空纤维函数 $F(x,\cdot)$, 它的所有非空水平集衰退锥也都为零; 存在性定理使其最优值有限并取到, 极小点集合紧. 因而 $p$ 不取 $-\infty$, 又在 $x_0$ 有限, 所以真, 上图闭则说明它闭.
\end{proof}
\begin{example}
原第69页反例取 $F(x,z)=e^{-\sqrt{xz}}$ 当 $x,z\ge0$, 其余处取 $+\infty$. 求 $p(x)=\inf_zF(x,z)$ 并解释其闭性.
\end{example}
\begin{solution}
几何平均 $\sqrt{xz}$ 在非负象限上凹, 所以 $-\sqrt{xz}$ 凸; 指数函数凸且单调递增, 复合仍凸. 非负象限闭, 函数在其中连续, 扩展后闭. 当 $x>0$, 令 $z\to\infty$ 得下确界0但不取到; 当 $x=0$, 值恒为1; 当 $x<0$, 无可行 $z$. 因而 $p(x)=0$ 对 $x>0$, $p(0)=1$, $p(x)=+\infty$ 对 $x<0$. 正数列趋于0使下半连续性失败. 此处每个非空纤维水平集都缺少定理要求的紧性.
\end{solution}

\originalsection{分离超平面}
\begin{definition}
超平面 $H=\{x:a\T x=b\}$ 要求 $a\ne0$. 若 $a\T x_1\le b\le a\T x_2$ 对所有 $x_i\in C_i$ 成立, 称 $H$ 分离两集. 若两边都为严格不等式, 称严格分离. 若分离且两集不全包含在 $H$ 中, 称适当分离. 若 $x\in\cl C$ 且经过 $x$ 的超平面把 $C$ 放在一侧, 称其在 $x$ 支持 $C$.
\end{definition}
\begin{theorem}\label{thm:separation}
非空凸集 $C$ 的非内点 $x$ 可与 $C$ 由经过 $x$ 的超平面分离. 两个不交非空凸集可分离. 若其中一个闭而另一个紧, 则可严格分离. 两个非空凸集可适当分离, 当且仅当它们的相对内部不交.
\end{theorem}
\begin{proof}
先对闭凸集及外点 $x$ 使用投影 $p=P_C(x)$. 向量 $a=p-x\ne0$ 满足 $a\T(y-p)\ge0$ 对 $y\in C$, 并且 $a\T p>a\T x$, 所以给出严格分离. 对边界点, 取外点 $x_k\to x$, 将各分离法向量规范化为单位向量并取收敛子序列, 极限仍为非零单位向量, 使 $a\T(y-x)\ge0$. 非闭集以闭包代替. 若仿射包是低维, 可取垂直于仿射包的法向量; 若 $x$ 不在闭包, 使用严格分离后把平面移至 $x$ 即可.

对不交集合, 在凸差集 $C_2-C_1$ 与原点间用第一条, 得 $a\T(x_2-x_1)\ge0$. 因而存在有限 $b$ 夹在 $\sup_{C_1}a\T x$ 与 $\inf_{C_2}a\T x$ 之间. 闭集与紧集的差闭且不含0; 投影原点到差集, 得到严格正的间隔, 可将平面放在间隔中点.

若相对内部不交, 则 $0\notin\ri(C_2-C_1)$, 因为相对内部与和交换. 在差集的仿射包内应用支持构造, 或当原点不在该仿射包时严格分离, 可选一个不在整个差集上恒为0的分离泛函, 因而适当分离成立. 反向若两集有共同相对内点 $z$, 任何分离泛函在该点取得两侧等号. 因为它在每个集合中都在相对内部取线性极值, 必须在两集各自恒定, 所以只能把两集都包含于平面中, 无法适当分离.
\end{proof}
适当分离在差集的仿射包内构造的细节如下. 若0位于相对边界, 外点可在仿射包内选择, 单位法向量也在其方向空间内; 极限法向量非零, 而整个差集仿射张成该空间, 所以它不能在差集上恒为0. 这避免了只取环境空间中无信息的垂直法向量.

\begin{corollary}
$\cl(\conv C)$ 等于所有包含 $C$ 的闭半空间的交.
\end{corollary}
\begin{proof}
每个包含 $C$ 的闭半空间都包含其闭凸包. 若 $x$ 在闭凸包外, 点与闭凸包的严格分离提供一个包含该包而排除 $x$ 的闭半空间, 因而反向包含成立.
\end{proof}

\originalsection{多面体分离的加强}
后续线性约束的对偶性不要求所有约束严格可行, 其基础是多面体分离可以加强到“不把非多面体一侧整个放在平面上”. 为说明有限性在哪里起作用, 先建立多面体的有限生成表示.

\begin{lemma}\label{lem:poly-finite}
每个多面体锥都由有限个向量生成. 每个非空多面体都可写成有限点的凸包与有限向量生成锥之和.
\end{lemma}
\begin{proof}
全空间由 $\pm e_i$ 有限生成. 若 $K=\cone\{v_i\}$, 则与半空间 $a\T x\le0$ 相交后, 可用满足 $a\T v_i\le0$ 的原生成向量及每一正负配对向量
\[
(-a\T v_j)v_i+(a\T v_i)v_j\quad(a\T v_i>0, a\T v_j<0)
\]
生成. 为验证, 将任意非负组合中正的 $a$ 分量与负的 $a$ 分量逐量配对抵消, 每次抵消正对应上述配对向量, 剩余只含非正分量. 有限次加入定义锥的半空间, 即得第一条; 等式用两个半空间表示.

对多面体 $P=\{x:Bx\le b,Ex=e\}$, 提升为锥 $K=\{(x,t):Bx\le bt,Ex=et,t\ge0\}$. 将有限生成向量中 $t>0$ 的项按其最后坐标归一到1, 其横坐标形成有限点集; $t=0$ 的项形成有限方向集. 切取 $t=1$ 时, 正高度项的系数和为1, 零高度项仍为非负组合, 得所述表示.
\end{proof}
\begin{theorem}\label{thm:polyseparation}
若 $C$ 非空凸, $P$ 非空多面体, 且 $\ri C\cap P=\varnothing$, 则存在分离超平面, 并且 $C$ 不全包含在该平面中.
\end{theorem}
\begin{proof}
先适当分离 $C,P$. 若平面不包含整个 $C$, 已得结论. 否则 $C$ 在平面上, 而 $P$ 在一侧且有点严格位于该侧. 将 $P$ 限制到与平面的交, 这是 $P$ 的真面, 维数下降; 若交为空, 多面体的有限生成点在该泛函上有统一正余量, 可将平面朝多面体方向小幅平移, 使它严格分离并不包含 $C$. 对剩余面与 $C$ 再作适当分离, 必要时重复. 每次若仍包含整个 $C$, 就再取一个真面. 维数有限, 不可能无限下降, 因而最终得到一个不包含整个 $C$ 的分离泛函.

还须把最后泛函扩展为对原 $P$ 的分离. 在引理\ref{lem:poly-finite}的有限点与射线表示上, 前一个泛函对不属于所取面的生成点或方向有严格正余量; 对面的生成项, 新泛函已经非负. 因而把新泛函乘以足够小的正数再加到前一个泛函上, 对全部有限生成项仍非负. 在 $C$ 上前一个泛函恒定, 所以加入新泛函后保留其非恒定性. 逐级反向合并, 就得到对原 $P$ 的分离且不包含整个 $C$ 的平面.
\end{proof}
有限生成项是此证明不可缺少的条件. 对非多面体, 无限生成方向的严格余量可能趋于0, 无法选取统一的小系数.

\originalsection{仿射下界与共轭函数}
\begin{lemma}\label{lem:affine-minorant}
每个真凸函数都具有有限的仿射下界. 若 $f$ 闭真凸且 $r<f(x)$, 则存在仿射函数 $\ell$ 满足 $\ell\le f$ 且 $\ell(x)>r$.
\end{lemma}
\begin{proof}
对真凸函数, 取 $x_0\in\ri\dom f$. 它在相对小球上有限连续. 在上图相对于 $\aff(\dom f)\times\R$ 的空间内, 分离上图和 $(x_0,f(x_0)-1)$. 分离泛函写成 $a\T x+\beta w\ge c$, 由于上图可向上无限移动, $\beta\ge0$. 若 $\beta=0$, 相对内点 $x_0$ 处的严格分离不可能成立, 所以可取得 $\beta>0$. 除以 $\beta$ 得仿射下界, 并把横向泛函从仿射包延拓到全空间. 这也完成上一章闭包真性的前向引用.

对于闭真凸函数, 点 $(x,r)$ 在闭上图外, 严格分离得到 $(a,\beta)$, $\beta\ge0$. 若 $\beta>0$ 直接得到所需仿射函数. 若 $\beta=0$, 将该严格分离不等式加上已有仿射下界对应不等式的小正倍数. 严格间隔在小扰动下仍为正, 竖直系数变为正, 除以它便得到 $\ell(x)>r$.
\end{proof}
\begin{definition}
函数 $f$ 的凸共轭为 $f^*(y)=\sup_x\{y\T x-f(x)\}$. 它记录所有斜率为 $y$ 的仿射下界中最大的截距: $y\T x-f^*(y)\le f(x)$. 二次共轭记为 $f^{**}$.
\end{definition}
共轭是仿射函数族的上确界, 所以始终闭凸, 即使 $f$ 本身不凸. 对真凸函数, 仿射下界说明其共轭至少在一个点有限, 因而共轭真; 又因原函数在某点有限, 共轭从不取 $-\infty$.

\begin{theorem}[共轭定理]\label{thm:biconjugate}
总有 $f^{**}\le f$. 若 $f$ 闭真凸, 则 $f^{**}=f$. 更一般地, 若 $f$ 的闭凸包函数真, 则 $f^{**}$ 等于这一闭凸包函数.
\end{theorem}
\begin{proof}
由 $f^*(y)\ge y\T x-f(x)$ 得 $y\T x-f^*(y)\le f(x)$, 对 $y$ 取上确界得第一条. 若 $f$ 闭真凸, 对任意实数 $r<f(x)$, 引理\ref{lem:affine-minorant}给出 $a\T z+b\le f(z)$ 且 $a\T x+b>r$. 由共轭定义, $f^*(a)\le-b$, 所以 $f^{**}(x)\ge a\T x-f^*(a)\ge a\T x+b>r$. 对所有 $r<f(x)$ 成立, 包括 $f(x)=+\infty$ 的情形, 因而 $f^{**}\ge f$.

共轭的定义只取线性泛函在上图上的下界, 对上图取凸包或闭包不改变这些下界. 因而 $f$ 与其闭凸包有相同共轭; 对真闭凸包使用已证明结论即可.
\end{proof}
非真性不能忽略. 令 $f(x)=-\infty$ 对 $x\le0$, $f(x)=+\infty$ 对 $x>0$, 它闭凸, 但 $f^*\equiv+\infty$, $f^{**}\equiv-\infty$, 不等于 $f$.

\begin{example}
\label{ex:conjugate-models}
计算仿射函数、绝对值、幂函数及正定二次函数的共轭.
\end{example}
\begin{solution}
对 $f(x)=a\T x-b$, 有 $f^*(y)=b$ 若 $y=a$, 否则 $+\infty$. 因为 $y\ne a$ 时沿 $y-a$ 方向的线性表达式无界. 对 $f(x)=|x|$, 当 $|y|\le1$ 时 $yx-|x|\le0$ 且0处取到; 当 $|y|>1$ 时沿对应符号方向无界, 所以共轭是区间 $[-1,1]$ 的指示函数.

取 $1<p<\infty$, $q=p/(p-1)$. 对 $f(x)=\sum_i|x_i|^p/p$, 各坐标独立, 最大化条件为 $y_i=\operatorname{sgn}(x_i)|x_i|^{p-1}$, 代回得到 $f^*(y)=\sum_i|y_i|^q/q$.

对 $f(x)=\tfrac12x\T Qx+a\T x+b$, $Q\succ0$, 最大化的一阶条件给出 $x=Q^{-1}(y-a)$, 完成平方得 $f^*(y)=\tfrac12(y-a)\T Q^{-1}(y-a)-b$. 一维 $f(x)=cx^2/2$ 的共轭因此为 $y^2/(2c)$.
\end{solution}
若 $f(x)=p(A(x-c))+a\T x+b$ 且 $A$ 可逆, 变量代换 $z=A(x-c)$ 给出
$f^*(y)=p^*(A^{-\mathsf T}(y-a))+c\T(y-a)-b$. 可逆条件是这里直接代换公式的一部分, 不可用于一般降维映射.

\originalsection{支撑函数、极锥与 Farkas 引理}
集合 $X$ 的支撑函数为 $\sigma_X(y)=\sup_{x\in X}y\T x$, 它是指示函数的共轭. $X,\cl X,\conv X,\cl\conv X$ 有相同支撑函数; 它的上图是闭凸锥. 支撑值是投影方向上的上确界, 可能为无穷或未取到, 因而不能总画成一个实际“最远点”.

本书使用极锥约定 $C^\circ=\{y:y\T x\le0\ \forall x\in C\}$. 在锥规划中将另用非负约定的对偶锥 $C^+=-C^\circ$, 两者的符号不可混用. 若 $C$ 是包含0的锥, 则 $\sigma_C=\ind{C^\circ}$, 因为任何正内积都可按比例放大至无穷. 共轭定理给出双极公式 $(C^\circ)^\circ=\cl\conv C$.

\begin{theorem}[Farkas 引理]\label{thm:farkas}
设 $A$ 的列为 $a_1,\ldots,a_m$. 方程 $A\lambda=b$, $\lambda\ge0$ 可解, 当且仅当对每个 $y$, $A\T y\le0$ 都蕴含 $b\T y\le0$. 等价地, 下列两者恰有一个成立: $b$ 是各列的非负组合; 存在 $y$ 使 $A\T y\le0$ 而 $b\T y>0$.
\end{theorem}
\begin{proof}
必要性由 $b\T y=\lambda\T A\T y\le0$. 对充分性, 列向量生成的锥闭, 已在 Carathéodory 推论中证明. 若 $b$ 不在该锥中, 对点与闭凸锥严格分离; 锥的齐次性强制分离泛函在整个锥上非正, 而在 $b$ 上严格正, 正是第二种条件. 两种条件不可能同时成立.
\end{proof}
这把“找一组非负系数”和“找一个证实不可能的线性泛函”连接起来. 在线性规划中, 前者给出对偶乘子, 后者给出降低目标的可行方向; 这一几何关系将在下一章展开.
